\chapter{Planetary boundary layer: YSU}\label{ch:pbl}

The planetary boundary layer (PBL) is the turbulent layer next to the ground, typically a few hundred metres deep at
night and one to three kilometres on a sunny afternoon. At 900~m grid spacing (d01), its eddies are not resolved, so
a scheme must mix heat, moisture and momentum vertically. The case uses the Yonsei University scheme, YSU
\citep{hong2006ysu}, on d01 (\code{bl_pbl_physics = 1}). On d02 (100~m) there is no PBL scheme: the large-eddy closure
of \cref{ch:turb} resolves the large eddies. The core is \src{phys/physics_mmm/bl_ysu.F90}{bl_ysu_run}, wrapped by
\code{phys/module_bl_ysu.F} and called from \code{phys/module_pbl_driver.F}.

\section{A non-local K profile}

Inside the mixed layer, YSU specifies the eddy diffusivity as a profile scaled by the boundary-layer height $h$ and a
velocity scale $w_s$:
\begin{equation}
  K_m(z) = \kappa\,w_s\,z\left(1 - \frac{z}{h}\right)^{p}, \qquad p = 2 \ (\code{pfac}).
\end{equation}
The flux of a scalar $C$ then has a local part and a non-local, \emph{counter-gradient} part that carries heat
upward even where the mean gradient is neutral or slightly stable:
\begin{equation}
  \overline{w'C'} = -K_c\left(\pd{C}{z} - \gamma_c\right) + \overline{(w'C')}_h\left(\frac{z}{h}\right)^3 .
\end{equation}
\begin{itemize}
  \item $\gamma_c$ is proportional to the surface flux divided by $w_s h$, with the coefficient $b = 6.8$
    (\code{bfac}).
  \item The last term is the \emph{explicit entrainment} flux at the top of the layer, the distinctive feature of YSU:
    the downward heat flux at $h$ is set as a fraction of the surface buoyancy flux, rather than left to the
    diffusivity profile.
\end{itemize}
The Prandtl number that turns $K_m$ into $K_h$ is bounded between 0.25 and 4 (\code{prmin}, \code{prmax}).

\section{Boundary-layer height}

$h$ is the lowest level where the bulk Richardson number, computed from the surface to that level,
\begin{equation}
  \mathrm{Ri}_B(z) = \frac{g\,(\theta_v(z) - \theta_s)\,z}{\theta_{va}\,\bigl(U(z)^2 + V(z)^2\bigr)},
\end{equation}
exceeds a critical value: 0 for unstable conditions (\code{brcr_ub}) and 0.25 for stable ones (\code{brcr_sb}). In
unstable conditions, $\theta_s$ includes a thermal excess proportional to the surface heat flux. The search walks up
the column level by level and interpolates between the two levels around the critical value.

\section{Above the boundary layer}

In the free atmosphere, YSU uses a local closure:
\begin{equation}
  K = l^2\,\left|\pd{\mathbf{V}}{z}\right|\,f(\mathrm{Ri}), \qquad \frac{1}{l} = \frac{1}{\kappa z} + \frac{1}{\lambda_0},
  \quad \lambda_0 = 30\ \text{m} \ (\code{rlam}),
\end{equation}
with a stability function of the local gradient Richardson number, and lower bounds of $0.1$ and
$0.01\ \text{m}^2\,\text{s}^{-1}$ for momentum and heat (\code{xkzminm}, \code{xkzminh}).

\section{The implicit vertical diffusion}

With the diffusivities known, the tendencies are computed by an implicit (backward) diffusion step, so that large $K$
in thin layers does not limit the time step. Each scalar gives one tridiagonal system per column
(\src{phys/physics_mmm/bl_ysu.F90}{tridin_ysu}). The two horizontal wind components are solved together with the
surface stress (\code{tridi2n}). The scheme returns tendencies of $u$, $v$, $\theta$, $q_v$, $q_c$ and $q_i$
(\code{RUBLTEN}, \dots), which the dynamics adds every step (\cref{ch:phys}).

\begin{algorithm}[ht]
\caption{One YSU column (simplified).}
\label{alg:ysu}
\begin{algorithmic}[1]
\State virtual potential temperature, wind speed and heights of the column
\State boundary-layer height from $\mathrm{Ri}_B$, first without and then with the thermal excess
\State velocity scale $w_s$, counter-gradient terms, entrainment fluxes at $h$
\State diffusivities: profile below $h$, local closure above
\State implicit diffusion of scalars (\code{tridin_ysu}) and of momentum (\code{tridi2n})
\State tendencies from the new minus the old profiles
\end{algorithmic}
\end{algorithm}

\begin{portnote}
\begin{itemize}
  \item \textbf{Out-of-bounds read.} The wrapper copies thirteen arrays used only with the urban BEP model. When the
    urban model is off, they are allocated with size one, so the copy reads out of bounds. The reference build guards
    the copy with the urban flag. The values are never used, so the guard changes no result.
  \item \textbf{Termination.} \code{get_pblh} contains a \code{DO WHILE} loop with no bound on its index. It terminates
    for physical profiles, and the column harness T-YSU-COL checks that it does on real data.
  \item \textbf{Structure.} Like WSM6, YSU becomes one GPU thread per column, with fixed-size temporaries (about fifty
    two-dimensional and forty one-dimensional arrays) and its tridiagonal solvers as device routines.
\end{itemize}
\end{portnote}

\section{Large-eddy simulation on d02}

At 100~m, the largest eddies of a convective boundary layer (about the depth of the layer) are resolved by about ten
to twenty grid points. Running a PBL scheme there would mix the same eddies twice. The case therefore switches it off
on d02 (\code{bl_pbl_physics = 0}) and uses the three-dimensional TKE closure of \cref{ch:turb}. The surface fluxes
of the surface layer and Noah enter as the lower boundary condition of the vertical diffusion.
